\subsection{Open-World Execution and Recovery}
\label{sec:rw_execution}

Even when perception and planning succeed, action execution in open-world environments may fail due to unexpected dynamics, novel object properties, or environmental changes. Recovery from such failures requires mechanisms that go beyond standard error handling: the agent must identify why the failure occurred and adapt its behavior accordingly.

Classical approaches to execution-time adaptation include plan monitoring and repair \citep{fox2006plan, fritz2007monitoring}, which detect deviations from expected execution outcomes and revise plans within the existing domain model. However, plan repair assumes that the domain model remains valid and that a correct plan exists within the current model, which may not hold when qualitative novelty has occurred. Continual learning methods \citep{kirkpatrick2017overcoming, rusu2016progressive} address the problem of acquiring new capabilities without forgetting previous ones, but most assume known task boundaries and distributions. Sim-to-real transfer techniques, including domain randomization \citep{tobin2017domain} and system identification \citep{peng2018sim}, improve policy robustness by training across varied conditions, but the variations are typically parameterized in advance and may not cover the kinds of qualitative changes encountered in open-world settings.

More recently, work within the DARPA SAIL-ON program \citep{senator2019science} has specifically targeted the problem of agents that must detect, characterize, and accommodate novelty in their environments. This program motivated research on novelty as an agent-relative concept: what constitutes novelty depends on the agent's prior knowledge and capabilities, not solely on properties of the environment \citep{boult2021towards, langley2020open}. Several approaches have explored how agents can adapt their policies or models in response to detected novelty \citep{balloch2022novgrid, goel2022rapidlearn}, though systematic methods for efficient recovery from qualitatively new failure modes remain limited.
